\begin{table}[t]
\centering
\caption{Experiment matrix and evidence coverage.}
\label{tab_experiment_matrix}
\scriptsize
\setlength{\tabcolsep}{3.5pt}
\renewcommand{\arraystretch}{1.08}
\begin{tabular}{@{}>{\raggedright\arraybackslash}p{0.10\textwidth}>{\raggedright\arraybackslash}p{0.20\textwidth}>{\raggedright\arraybackslash}p{0.29\textwidth}>{\raggedright\arraybackslash}p{0.32\textwidth}@{}}
\toprule
Block & Purpose & Methods & Evidence coverage \\
\midrule
P1 & Core and artifact diagnostics & PPO backbone, Risk-only, No-action guard, Guard, Shield, Gated-risk & 4 densities; 24 aggregate cells; raw and per-seed CSVs \\
P2 & External baselines & PPO backbone, RSS/TTC filter, PPO-Lagrangian, Guard, Shield, Gated-risk & density 0.15; 720 diagnostic episodes represented in aggregate rows \\
P4 & Density stress & PPO backbone plus runtime variants and artifact controls & densities 0.08, 0.15, 0.20, 0.25 \\
P5 & Sensitivity & Guard, Shield, Gated-risk & target-speed sweep and TTC-threshold sweep; 21 aggregate cells \\
CI/seed & Uncertainty and reproducibility & all formal density-0.15 methods & Wilson intervals, bootstrap intervals, seed scatter, runtime by seed \\
\bottomrule
\end{tabular}
\begin{tablenotes}
\item The manuscript reports aggregate tables, while raw episode-level CSVs remain in the evidence package.
\end{tablenotes}
\end{table}
